\providecommand{\CoinPkFair}{}\renewcommand{\CoinPkFair}{0.04395}
\providecommand{\CoinPkBias}{}\renewcommand{\CoinPkBias}{0.302}
\providecommand{\CoinSimFair}{}\renewcommand{\CoinSimFair}{0.04105}
\providecommand{\CoinSimBias}{}\renewcommand{\CoinSimBias}{0.3024}
\providecommand{\CoinSimErr}{}\renewcommand{\CoinSimErr}{0.001449}
\providecommand{\CoinPostBias}{}\renewcommand{\CoinPostBias}{0.873}
\providecommand{\CoinPostBiasFive}{}\renewcommand{\CoinPostBiasFive}{0.09696}
\providecommand{\CoinLR}{}\renewcommand{\CoinLR}{6.872}
\providecommand{\CoinNexp}{}\renewcommand{\CoinNexp}{20000}
\providecommand{\zeroaShrinkFig}[6]{%
\begin{tikzpicture}[scale=0.8]
  \pgfmathsetmacro{\zA}{6*#1}                 %
  \pgfmathsetmacro{\zL}{\zA-6*#1*#2}           %
  \pgfmathsetmacro{\zR}{\zA+6*(1-#1)*#3}       %
  \pgfmathsetmacro{\zJH}{#1*#2}                %
  \pgfmathsetmacro{\zJN}{(1-#1)*#3}            %
  \pgfmathsetmacro{\zPE}{\zJH+\zJN}            %
  \pgfmathsetmacro{\zPost}{\zJH/\zPE}
  \pgfmathsetmacro{\zPostN}{1-\zPost}
  \pgfmathsetmacro{\zPH}{#1}
  \pgfmathsetmacro{\zPN}{1-#1}
  \begin{scope}
    \fill[formblue!18] (0,0) rectangle (6,4);
    \fill[tangentgreen!45] (\zL,0) rectangle (\zR,4);
    \draw (\zA,0) -- (\zA,4);
    \draw[thick] (0,0) rectangle (6,4);
    \draw[line width=1.6pt,fibreorange] (\zL,0) rectangle (\zR,4);
    \node[lbl] at ({\zL/2},2) {#4};
    \node[lbl] at ({(\zR+6)/2},2) {#5};
    \draw[faint] ({(\zL+\zA)/2},1.0) -- ({\zL-0.3},0.75)
      node[lbl,font=\footnotesize,anchor=east,black,inner sep=1pt]
      {\pgfmathprintnumber[fixed,precision=3]{\zJH}};
    \draw[faint] (\zR,2.6) -- ({\zR+0.45},2.95)
      node[lbl,font=\footnotesize,anchor=south west,black,inner sep=1pt]
      {\pgfmathprintnumber[fixed,precision=3]{\zJN}};
    \node[lbl,fibreorange!80!black,anchor=south] at ({(\zL+\zR)/2},4.05) {$E$: #6};
    \node[lbl] at (6.25,3.75) {$\Omega$};
    \draw[faint] (0,-0.15) -- (0,-0.3) -- (\zA,-0.3) -- (\zA,-0.15);
    \draw[faint] (\zA,-0.15) -- (\zA,-0.3) -- (6,-0.3) -- (6,-0.15);
    \node[lbl] at ({\zA/2},-0.6) {\pgfmathprintnumber{\zPH}};
    \node[lbl] at ({(\zA+6)/2},-0.6) {\pgfmathprintnumber{\zPN}};
    \node[lbl,align=center,text width=6.2cm,anchor=north] at (3,-0.95)
      {before: one strip for each coin,\\ as wide as its prior};
  \end{scope}
  \draw[->,thick,accent] (6.7,2) -- (8.1,2);
  \node[lbl,accent,align=center] at (7.4,2.75) {observe\\$E$};
  \begin{scope}[xshift=8.7cm]
    \draw[faint,dashed] (0,0) rectangle (6,4);
    \draw[faint,dashed] (\zA,0) -- (\zA,4);
    \fill[tangentgreen!45] (\zL,0) rectangle (\zR,4);
    \draw[very thin,black!40] (\zA,0) -- (\zA,4);
    \draw[line width=1.6pt,fibreorange] (\zL,0) rectangle (\zR,4);
    \draw[faint] ({(\zL+\zA)/2},2.4) -- ({\zL-0.25},2.9)
      node[lbl,anchor=east,black,align=right,inner sep=1pt]
      {#4\\ \pgfmathprintnumber[fixed,precision=3]{\zPost}};
    \draw[faint] ({\zR},1.2) -- ({\zR+0.45},0.9)
      node[lbl,anchor=west,black,align=left,inner sep=1pt]
      {#5\\ \pgfmathprintnumber[fixed,precision=3]{\zPostN}};
    \node[lbl,fibreorange!80!black,anchor=south] at ({(\zL+\zR)/2},4.05)
      {$E$, area \pgfmathprintnumber[fixed,precision=3]{\zPE}};
    \node[lbl,align=center,text width=6.2cm,anchor=north] at (3,-0.95)
      {after: $E$ is the new sample space,\\ each coin keeps its share of it};
  \end{scope}
\end{tikzpicture}}
 \chapter*{Preface}
\addcontentsline{toc}{chapter}{Preface}
\markboth{Preface}{Preface}
\setcounter{figure}{0}\setcounter{wdef}{0}\setcounter{wex}{0}\setcounter{lstlisting}{0}
\renewcommand{\thefigure}{P.\arabic{figure}}
\renewcommand{\thewdef}{P.\arabic{wdef}}
\renewcommand{\thewex}{P.\arabic{wex}}
\renewcommand{\thelstlisting}{P.\arabic{lstlisting}}

\section*{Who made these notes}
These notes were made by Chandra Prakash
(\texttt{chandra.pp@alumni.iitg.ac.in}; the latest version lives at
\href{https://lunar-photon.github.io}{\texttt{lunar-photon.github.io}})
together with Claude, an AI model made by Anthropic. The goals, the vision
and the way every idea is framed are Chandra's: which questions matter, which
pictures make them clear, and how the whole subject hangs together for a
physicist. The drafting of the text and the writing of the code were done
with Claude's help, one section at a time, with Chandra reading, correcting
and redirecting. It was a joint project and an ambitious one: to take a
physicist who has never properly learned statistics all the way from coin
tosses to Monte Carlo simulations of the microwave sky, Markov chain samplers,
Fisher forecasts and the topology of random fields, and to have every step
run as a Python script on the reader's own machine.

\section*{Goals and vision}
\paragraph{What the notes are for.}
The reader we have in mind is a phenomenologist working in particle physics,
cosmology or gravitational waves. You are at home with calculus, linear
algebra, Fourier analysis, quantum mechanics and statistical mechanics. You
use likelihoods, error bars and MCMC chains in your work, often through code
written by someone else, and you would like to know what each tool really
does and why it exists. You also want to \emph{implement} it: to write the
fifty lines of Python that do the job, see them work on data, and then trust
the standard library because you know what it is doing inside.

\paragraph{The vision.}
Three convictions shape every chapter.
\begin{enumerate}
  \item \emph{Every idea opens with a story.} Before any definition we meet a
        concrete situation that makes the idea necessary, and a physical
        picture of what the idea will do for us. The formula comes after we
        know what question it answers.
  \item \emph{Probability runs in two directions.} Prediction goes forward:
        given a hypothesis, how often does each outcome occur? Inference goes
        backward: given the outcome we saw, which hypotheses could lie behind
        it, and how much does each deserve belief now? Bayes' theorem is the bridge.
        Almost every confusion in statistics is a confusion between these two
        directions.
  \item \emph{One pipeline organises the whole subject.} A physical model
        produces a field or a data set. We choose what information in it
        matters, compress the data into a summary $S$, work out how $S$
        fluctuates from one realisation of the noise to the next (its sampling
        distribution $p(S\given\params)$), and only then infer the physics
        (\cref{pre:fig:pipeline}). The book walks along this pipeline, from
        probability distributions to Bayes, Monte Carlo, MCMC and the Fisher
        matrix, and ends where the summaries become most interesting:
        topological summaries of random fields.
\end{enumerate}

\begin{figure}[htbp]
\centering
\begin{tikzpicture}[
  stage/.style={draw=accent,thick,rounded corners=2pt,fill=ideabg,align=center,
                font=\small,text width=2.35cm,minimum height=1.15cm},
  flow/.style={->,thick,accent}]
  \node[stage] (m) at (0,0)     {physical\\ model $\params$};
  \node[stage] (d) at (3.0,0)   {field or\\ data set $\data$};
  \node[stage] (c) at (6.0,0)   {what\\ information\\ matters?};
  \node[stage] (s) at (6.0,-2.1){summary\\ $S(\data)$};
  \node[stage] (p) at (3.0,-2.1){sampling\\ distribution\\ $p(S\given\params)$};
  \node[stage] (i) at (0,-2.1)  {inference\\ on $\params$};
  \draw[flow] (m) -- (d);
  \draw[flow] (d) -- (c);
  \draw[flow] (c) -- (s);
  \draw[flow] (s) -- (p);
  \draw[flow] (p) -- (i);
\end{tikzpicture}
\caption{The pipeline that organises the book. The top row is the physics
and the observation; the bottom row is the statistics. Each chapter serves
one or two of these steps, and \cref{ch:00} places every chapter on its step.}
\label{pre:fig:pipeline}
\end{figure}

\paragraph{What you will be able to do at the end.}
You will be able to look at a physics measurement and say which of the two
directions each part of the analysis runs in. You will be able to write down
a likelihood for Poisson counts in a collider bin, for a Gaussian random field
on the sphere and for a gravitational-wave detector with coloured noise, and
explain why these are the same idea. You will be able to write a Monte Carlo
simulation that checks an analytic result, and to recognise when the same
simulation is no longer a check but the only instrument available. You will
be able to code a Metropolis--Hastings sampler from scratch and diagnose it,
compute a Fisher forecast and know when it can be trusted, and treat a
topological summary of a field (its Betti numbers, its Euler characteristic,
its persistence diagram) as one more random observable, with its own
covariance and likelihood.

\section*{How the book is organised}
\Cref{ch:00} is the map: the two directions of probability, the list of
questions the book answers, the pipeline, the menu of summaries, and how to
run the code. \Cref{ch:T0} is the small part of Python that every script in
the book uses, and how to tell whether a number a script prints can be
trusted.

The first part builds the statistics. \Cref{ch:01} and \cref{ch:02} build
probability and the distributions a physicist meets. \Cref{ch:03} turns data
into estimates and introduces the likelihood; \cref{ch:04} asks whether a
signal is really there and how to set a limit when it is not; \cref{ch:05} is
Bayesian inference.

The second part collects the physics the statistics is used on, one toolkit
per field, each a map of the objects and equations that the later chapters
need: cosmology and weak lensing (\cref{ch:T1}), astrophysics and
gravitational waves (\cref{ch:T2}), the morphology of the microwave sky
(\cref{ch:T3}), and collider particle physics (\cref{ch:T4}). They can be
read in order or consulted when a chapter points to them.

The third part turns the statistics into instruments. \Cref{ch:06} is Monte
Carlo and the pseudo-random numbers behind it. \Cref{ch:07} makes the
summary concrete for fields ($\xi(r)$, $P(k)$, $\Cell$) and \cref{ch:08} is
a Monte Carlo research problem on the microwave sky. \Cref{ch:09} and
\cref{ch:10} turn sampling distributions into parameter constraints with
MCMC and the Fisher matrix, and \cref{ch:11} applies all of it to galaxy
surveys and the local Universe.

The last part goes beyond two-point statistics and onto real data.
\Cref{ch:12} is topological data analysis of random fields, \cref{ch:13}
runs four complete analyses on public data (the Planck sky, its topology,
SDSS galaxies and a Higgs search), and \cref{ch:14} pulls the threads
together: given a new problem, which idea to reach for. Every chapter opens
with a roadmap and closes with a summary table.

\section*{Bayes' theorem in one picture, told with one box of each kind}
The text of each chapter carries the whole argument, and it reads correctly
from start to finish if every box is skipped. Boxes are a sparse sidebar:
something to find again, or to skip on a first reading. Rather than describe
them, we show one real box of each kind. All of them tell one story, the
story that also opens \cref{ch:00}: Bayes' theorem as information shrinking
the sample space.

\begin{story}[the coin from the drawer]
A drawer holds two coins that look identical. One is fair. The other is
bent and lands heads 80\% of the time. A friend picks one at random, tosses
it ten times, and reports 8 heads. Which coin was it?

\smallskip\noindent{\footnotesize\textsf{Stories give the concrete situation that makes an
idea necessary. In the chapters they are ordinary prose, the first paragraph
of a section.}}
\end{story}

The story contains both directions at once. If we knew which coin it was, we
could predict how often it gives 8 heads. We do not know the coin; we know the
8 heads. So the question runs backward, from the outcome to the coin.

\begin{framing}
A question about chance can face either way. Facing forward, we fix the
situation and ask which outcomes it can give and how often each one turns up.
Facing backward, the outcome is already in hand, and we ask which of the
situations that might have been true could have given it, and how much each
of them deserves to be believed now that we have seen it.

For the coin, prediction asks $\Prob(E\given H)$: \emph{given this coin, how likely was
this outcome?} Inference asks $\Prob(H\given E)$: \emph{given this outcome,
which way could it have happened, and how probable is each way?} The two
numbers are different, and Bayes' theorem converts one into the other.

\smallskip\noindent{\footnotesize\textsf{The central framing states the question a chapter
answers.}}
\end{framing}

To make the conversion we need one definition. We will look it up many times,
so it is worth having in a box.

\begin{workdef}[conditional probability]
For events $A$ and $B$ with $\Prob(A)>0$, the \newterm{conditional
probability} of $B$ given $A$ is
\[
  \Prob(B\given A)=\frac{\Prob(A\cap B)}{\Prob(A)} .
\]
It is the probability of $B$ once we know that $A$ happened
\unpack{we keep only the outcomes in which $A$ is true} and
renormalise \unpack{divide by $\Prob(A)$, so the kept outcomes again have
total probability 1}. Operationally: in a long run of repetitions, throw away
every run in which $A$ failed and count the fraction of the survivors in
which $B$ happened.

\smallskip\noindent{\footnotesize\textsf{Working definitions keep the precise sentence and
unpack every term a forgetful reader may need, in green brackets.}}
\end{workdef}

Now the reasoning, in words, before any number.

\begin{reasoning}
The backward question can always be said in four plain steps before any
equation is written:
\begin{enumerate}
  \item What did we actually see?
  \item Which explanations could have given it?
  \item If each explanation were true, how probable would that sight have been?
  \item How do those answers change the weight each explanation had before?
\end{enumerate}
For the coin: (1) $E$ = ``8 heads in 10 tosses''; (2) $B$ (bent coin) and $F$
(fair coin), each picked with probability $\tfrac12$; (3) the forward
question, answered by the binomial distribution; (4) keep only the paths that
end in $E$, and share the probability among them in proportion to how much
each contributes.

\smallskip\noindent{\footnotesize\textsf{Reasoning boxes state a question in words before
any symbol appears.}}
\end{reasoning}

Step 3 is ordinary prediction: if this coin were the one, how often would it
give 8 heads? For $n$ tosses of a coin with heads probability $p$, the
probability of exactly $k$ heads is $\binom nk p^k(1-p)^{n-k}$. The factor
$p^k(1-p)^{n-k}$ is the probability of one particular sequence with $k$ heads,
and $\binom nk$ counts how many such sequences there are. With $n=10$ and
$k=8$, the fair coin gives $\Prob(E\given F)=\binom{10}{8}/2^{10}=45/1024
=\CoinPkFair$ and the bent coin gives
$\Prob(E\given B)=45\times0.8^8\times0.2^2=\CoinPkBias$. Step 4 is Bayes'
theorem, and it has a picture.

\begin{intuition}[information shrinks the sample space]
Before the evidence, the whole sample space is relevant. Draw it as a
rectangle of area one and cut it into two strips, one for each coin, each as
wide as its prior, $\Prob(B)=\Prob(F)=\tfrac12$. Inside each strip, mark the part
in which 8 heads happens. It fills the fraction $\Prob(E\given B)=\CoinPkBias$
of the bent-coin strip and only $\Prob(E\given F)=\CoinPkFair$ of the fair-coin
strip, so the two pieces have areas $\Prob(E\cap B)$ and $\Prob(E\cap F)$.
Together they are the event $E$, one box straddling the line between the
strips. The evidence tells us that we are somewhere inside that box, so
everything outside it is ruled out: the sample space has shrunk to the box.
The posterior $\Prob(B\given E)$ is the share of the box that the bent-coin
piece fills (\cref{pre:fig:shrink}).

\smallskip\noindent{\footnotesize\textsf{Intuition boxes give the physical or visual
picture.}}
\end{intuition}

\begin{figure}[htbp]
\centering
\zeroaShrinkFig{0.5}{\CoinPkBias}{\CoinPkFair}{bent}{fair}{8 heads in 10}
\caption{Bayes' theorem as information shrinking the sample space, drawn to
scale. Left: the rectangle $\Omega$ has area 1, and since the coin is either
bent or fair, the two hypotheses fill it in two strips of width $\tfrac12$. The
event $E$, ``8 heads in 10'' (orange outline), covers the fraction $\CoinPkBias$
of the bent-coin strip, because that is how often the bent coin gives 8 heads,
and only $\CoinPkFair$ of the fair-coin strip. Each green piece therefore has
area prior times likelihood. Right: hearing ``8 heads'' discards everything
outside the orange box, and each coin keeps the share of the box that its piece
fills: $\CoinPostBias$ for the bent coin.}
\label{pre:fig:shrink}
\end{figure}

The picture turns into a calculation with no new idea in it.

\begin{example}[the posterior for the coin]
With $\Prob(B)=\Prob(F)=\tfrac12$,
\begin{align*}
  \Prob(B\given E)
  &\eqby{1}\frac{\Prob(E\cap B)}{\Prob(E\cap B)+\Prob(E\cap F)}\\
  &\eqby{2}\frac{\Prob(E\given B)\Prob(B)}{\Prob(E\given B)\Prob(B)+\Prob(E\given F)\Prob(F)}\\
  &\eqby{3}\frac{\Prob(E\given B)}{\Prob(E\given B)+\Prob(E\given F)}
   =\frac{\CoinPkBias}{\CoinPkBias+\CoinPkFair}=\CoinPostBias .
\end{align*}
\begin{explain}
  \item Conditional probability, $\Prob(B\given E)=\Prob(E\cap B)/\Prob(E)$,
        with $\Prob(E)$ split into the two disjoint ways $E$ can happen (the
        two green pieces of the box).
  \item The product rule $\Prob(E\cap H)=\Prob(E\given H)\Prob(H)$ for each
        piece: its share of the strip times the strip's width.
  \item The equal priors $\tfrac12$ cancel between numerator and denominator.
\end{explain}
The 8 heads moved the probability of the bent coin from $0.5$ to
$\CoinPostBias$.

\smallskip\noindent{\footnotesize\textsf{Examples are solved to the last line. Every
non-trivial equals sign is labelled and explained right after the chain.}}
\end{example}

One detail of the example is easy to overlook: the prior widths of the two
strips cancelled only because they were equal. With unequal priors they stay
in the answer.

\begin{takeaway}
In the picture, Bayes' theorem needs no extra assumption: it falls out of
the picture. There is one patch, the overlap $H\cap E$, and two ways to
measure it. As a share of the $H$ strip it is $\Prob(E\given H)$; as a share
of the box $E$ it is $\Prob(H\given E)$. Bayes' theorem is the
statement that both shares describe the same area,
$\Prob(E\given H)\Prob(H)=\Prob(H\given E)\Prob(E)$.

For the coin, with $H=B$: a forward probability is the \emph{share of a strip}
that the evidence covers; a joint probability is an \emph{area}; the posterior
is a \emph{ratio of areas} once the evidence has cut the sample space down to
the box. Move the line between the strips and both pieces change width, so the
same data give a different posterior.

\smallskip\noindent{\footnotesize\textsf{Takeaways say in a sentence or two what a
computation taught us.}}
\end{takeaway}

The most common error in reading such results is to swap the two directions,
and it has a famous statistical form.

\begin{pitfall}[a $p$-value is not the probability that the null is true]
Suppose we test ``the coin is fair'' by asking how often a fair coin gives 8
or more heads in 10 tosses. That number, $(45+10+1)/1024\approx0.055$, is a
forward probability, $\Prob(\text{this or more extreme}\given F)$. It is not
the probability that the coin is fair, $\Prob(F\given E)$. That we just found
to be $1-\CoinPostBias$ with equal priors, and it depends on the prior as
well. \Cref{ch:04} makes this precise.

\smallskip\noindent{\footnotesize\textsf{Pitfalls flag a genuine trap.}}
\end{pitfall}

Every step of the story is a few lines of code. The script behind the numbers
above computes the forward probabilities from the binomial formula, checks
them by simulating \CoinNexp{} ten-toss experiments with each coin, and then
runs Bayes on the observed outcome. The backward step is four lines.

\pyfile[firstline=42,lastline=45,firstnumber=42]{code/ch00/01_coin_two_directions.py}{The
backward direction: keep the paths that end at the observed $k$, then
normalise}

Line 43 is the area of the bent-coin piece of the box in \cref{pre:fig:shrink},
line 44 the area of the fair-coin piece, and line 45 divides the first by their
sum.
The arrays run over every possible $k$ at once, so the script produces the
posterior for all eleven outcomes in one go. For $k=8$ it prints
$\CoinPostBias$; for $k=5$ it gives $\CoinPostBiasFive$. Every number quoted
in the text is written to a file by the script that computed it and read back
into \LaTeX. So the text and the code cannot drift apart (\cref{ch:00}
describes the machinery).

The simulation in the same script is a small instance of an idea that runs
through the book: a Monte Carlo experiment can check a formula, or it can
replace one.

\begin{inpractice}[verification versus research]
Here the simulation \emph{verifies} a known answer. The fraction of
simulated fair-coin experiments with 8 heads is $\CoinSimFair$, against the
exact $\CoinPkFair$. The difference is about two Monte Carlo standard errors
($\pm\CoinSimErr$ each), an ordinary fluctuation that shrinks as $1/\sqrt{N}$
when more experiments $N$ are simulated. In research the same machinery is
often the only instrument. The covariance of a power-spectrum estimator on a
masked, noisy microwave sky has no closed form, and thousands of simulated
skies are how it is measured (\cref{ch:08}).

\smallskip\noindent{\footnotesize\textsf{These boxes separate a numerical check of known
theory from the same tool used where no analytic answer exists.}}
\end{inpractice}

When an idea has several tools, a small table says which to use when.

\begin{taxonomy}[the coin, asked four ways]
\small
\begin{tabular}{@{}p{0.42\linewidth}p{0.3\linewidth}l@{}}
  \textbf{question in words} & \textbf{object} & \textbf{chapter}\\
  how often does a fair coin give 8 heads? & $\Prob(k\given p)$ & \cref{ch:02}\\
  what is our best guess of $p$? & estimator $\hat p=k/n$ & \cref{ch:03}\\
  is 8 heads surprising for a fair coin? & $p$-value & \cref{ch:04}\\
  how probable is each value of $p$? & posterior $p(p\given k)$ & \cref{ch:05}\\
\end{tabular}

\smallskip\noindent{\footnotesize\textsf{Taxonomies list which tool answers which
question.}}
\end{taxonomy}

The picture is not ours alone; the literature draws it in
several ways.

\begin{sourcesays}[Wasserman p.~ix, Fig.~1]
Wasserman opens with one diagram: a data generating process and the observed
data, joined by an arrow labelled ``probability'' going forward and one
labelled ``inference'' coming back \mainref{p.~ix}. Kruschke describes the
backward arrow as a reallocation of credibility across possibilities
\srcref{Kruschke}{pp.~16--22}. In the chapters, such links are short
citations at the point of use, with one ``Reading the literature'' list per
chapter.

\smallskip\noindent{\footnotesize\textsf{Source boxes map the notes onto the books, and flag
their misprints.}}
\end{sourcesays}

\section*{The labelled-equality convention}
Every calculation is written as a chain. Each non-trivial equals sign carries
a small label, and the labels are explained, in order, right after the chain.
For the odds form of the coin problem:
\begin{align*}
  \frac{\Prob(B\given E)}{\Prob(F\given E)}
  &\eqby{1}\frac{\Prob(E\given B)\,\Prob(B)/\Prob(E)}{\Prob(E\given F)\,\Prob(F)/\Prob(E)}
  \eqby{2}\frac{\Prob(E\given B)}{\Prob(E\given F)}\cdot\frac{\Prob(B)}{\Prob(F)}
  =\CoinLR\times1 .
\end{align*}
\begin{explain}
  \item Bayes' theorem for each coin.
  \item $\Prob(E)$ cancels. Posterior odds are the likelihood ratio times the
        prior odds.
\end{explain}
Odds of $\CoinLR$ to 1 are a probability $\CoinLR/(1+\CoinLR)=\CoinPostBias$,
as before.

\section*{The code}
Every figure made from data and every simulated number comes from a short
script in \texttt{code/chNN/}, runnable from the Notes directory with
\texttt{python3 code/chNN/NN\_name.py}. Randomness is seeded from the script's
name, so rerunning a script reproduces its figures bit for bit. Each method
is first implemented from scratch, so the mechanism is visible, and then
checked against a standard library (\texttt{emcee}, \texttt{getdist},
\texttt{healpy}, \texttt{gudhi} and others). \Cref{ch:00} ends with a tour of
the layout and the shared helpers.

\section*{Sources}
\begin{itemize}
  \item L.~Wasserman, \emph{All of Statistics} (Springer 2004), followed
        chapter by chapter and cited as \mainref{p.~\dots} with printed page
        numbers.
  \item J.~K.~Kruschke, \emph{Doing Bayesian Data Analysis} (2nd ed.), the
        model of clarity, cited as \srcref{Kruschke}{\dots}.
  \item G.~Cowan, \emph{Statistical Data Analysis}; D.~S.~Sivia,
        \emph{Data Analysis: A Bayesian Tutorial}; G.~Casella and
        R.~L.~Berger, \emph{Statistical Inference}; B.~Lambert, \emph{A
        Student's Guide to Bayesian Statistics}; M.~P.~Hobson et al.,
        \emph{Bayesian Methods in Cosmology}; and \emph{Seeing Theory}, each
        cited by name where it is clearest.
  \item Research papers on cosmological inference, pseudo-$\Cell$ estimation,
        gravitational-wave forecasts and the topology of random fields, cited
        in the usual way and collected in the bibliography.
\end{itemize}

\begin{recap}[the dictionary of this preface]
\begin{tabular}{@{}ll@{}}
  prediction & $\Prob(E\given H)$, the share of the $H$ strip covered by $E$\\
  joint probability & $\Prob(E\cap H)=\Prob(E\given H)\Prob(H)$, an area\\
  evidence & $\Prob(E)$, the total surviving area\\
  inference & $\Prob(H\given E)$, a ratio of surviving areas\\
  odds form & posterior odds $=$ likelihood ratio $\times$ prior odds\\
  the pipeline & model $\to$ data $\to$ summary $\to p(S\given\params)\to$ inference\\
\end{tabular}

\smallskip\noindent{\footnotesize\textsf{A recap closes every part.}}
\end{recap}

\renewcommand{\thefigure}{\thechapter.\arabic{figure}}
\renewcommand{\thewdef}{\thechapter.\arabic{wdef}}
\renewcommand{\thewex}{\thechapter.\arabic{wex}}
\renewcommand{\thelstlisting}{\thechapter.\arabic{lstlisting}}
